\section{Discussion}
\label{sec:discussion}

\subsection{The recommendation}

For a build long enough that a developer disengages, deliver the result
ambiently and in place. A gutter marker, an underline, or a status change on
the file tab all satisfy the requirement, which is that the information be
available in the periphery of the region the developer is already looking at
and require no action to clear. This is not a small optimisation. Halving
resumption lag on a notification that fires several times an hour is worth more
than most latency work on the build itself, and it costs nothing to implement.

\subsection{When it stops holding}

The recommendation has a floor. We ran a follow-up with 12 of our participants
at a build duration of 6~s rather than 42~s, and the ordering inverts:
\emph{Ambient} produces a mean lag of 4.1~s against 2.8~s for
\emph{Immediate}. At six seconds developers do not disengage. They wait, gaze
fixed on the edit region, and an ambient marker that changes colour outside
foveal attention is simply missed for a second or two, whereas a toast is
noticed at once. The crossover in our follow-up sits between 8 and 12~s,
though with 12 participants we treat that interval as indicative rather than
established.

This gives a concrete design rule that a tool can implement without asking the
user anything: measure the build, and choose the modality from its duration.
Below roughly ten seconds, notify immediately. Above it, notify ambiently.

\subsection{Why batching is not the answer}

The batched condition is second best on every measure, and it is worse than
ambient on all four. Batching addresses timing, deferring the interruption to a
typing pause, but it preserves the modality: the toast still appears in the
corner, still demands a gaze shift, and still requires dismissal. Our gaze data
suggest the modality is the larger of the two costs, so a design that fixes
only the timing captures only part of the available benefit. This also predicts
that batched ambient markers, which we did not test, should be no better than
ambient markers alone, since the ambient marker imposes no cost worth deferring.

\subsection{Generalisation beyond compilation}

Nothing in the mechanism is specific to builds. Test runs, deployments, long
queries, and language server indexing all produce an asynchronous result that
arrives after the developer has moved on. We would expect the same ordering
wherever the wait exceeds the disengagement threshold and the result can be
rendered in place. Where the result cannot be rendered in place, because it
concerns a file the developer is not looking at, the ambient design loses its
advantage and the question becomes open again.
